\documentclass{article}

% Packages
\usepackage{microtype}
\usepackage{graphicx}
\usepackage{subcaption}
\usepackage{booktabs} 
\usepackage{hyperref}
\usepackage{amsmath}
\usepackage{amssymb}
\usepackage{mathtools}
\usepackage{amsthm}
\usepackage[capitalize,noabbrev]{cleveref}

% For pseudocode
\usepackage{algorithm}
\usepackage{algorithmic}

% ICML Style
\usepackage{icml2026}

% Running title
\icmltitlerunning{Mitigating Calibration Variance in Data-Efficient BatchNorm Calibration for Quantized Model Merging}

\begin{document}

\twocolumn[
\icmltitle{Mitigating Calibration Variance in Data-Efficient BatchNorm Calibration \\ for Quantized Model Merging}

\icmlsetsymbol{equal}{*}

\begin{icmlauthorlist}
\icmlauthor{The Empiricist}{equal,inst}
\end{icmlauthorlist}

\icmlaffiliation{inst}{Department of Empirical Machine Learning, Institute of Large-Scale Experimentation}
\icmlcorrespondingauthor{The Empiricist}{empiricist@experiments.org}

\icmlkeywords{Model Merging, Post-Training Quantization, BatchNorm Calibration, Calibration Variance, Robustness, Deep Learning}

\vskip 0.3in
]

\printAffiliationsAndNotice{}

\begin{abstract}
While multi-task model merging represents a highly storage-efficient, training-free paradigm for consolidating specialized expert networks, linear parameter combination often triggers severe performance collapse under aggressive post-training quantization (PTQ). Recent work has shown that a simple, task-specific, data-efficient BatchNorm calibration (DE-BN) using a tiny support set (e.g., 32 samples) can dramatically heal representation collapse. In this paper, we adopt an empirical, rigorous approach to expose a major, previously unaddressed vulnerability of this practice: \textbf{calibration variance}. We demonstrate that under low-bit (4-bit) quantization, DE-BN's performance is extremely sensitive to the random choice of calibration samples, leading to substantial variance in multi-task accuracy. To mitigate this, we propose two novel, simple, and highly effective protocols: \textbf{Robust Multi-Seed BatchNorm Calibration (RMS-BC)} and \textbf{Data-Efficient Momentum BatchNorm Calibration (DEM-BC)}. Through exhaustive empirical evaluation of ResNet-18 expert models merged across MNIST, Fashion-MNIST, and CIFAR-10, we show that RMS-BC reduces calibration variance by over 75\% while DEM-BC with a tuned momentum parameter $\beta = 0.1$ dramatically improves both average accuracy and stability, outperforming standard DE-BN by substantial absolute accuracy margins (up to +11.28\% in 4-bit regimes). Our work establishes a new, robust standard for data-efficient calibration in quantized model merging.
\end{abstract}

\section{Introduction}
\label{submission-introduction}
The development of specialized expert models fine-tuned from a shared pre-trained progenitor backbone has revolutionized downstream machine learning applications. In modern computer vision and natural language processing, practitioners frequently fine-tune a common base model—such as ResNet \cite{he16}, BERT \cite{devlin18}, or LLaMA \cite{touvron23}—on various distinct tasks, resulting in a suite of specialized expert networks. However, storing and serving dedicated high-capacity checkpoints for each downstream task introduces severe computational, memory, and logistical overheads, particularly in edge devices and resource-constrained environments. 

Multi-task model merging has emerged as an exceptionally appealing paradigm, enabling the zero-shot consolidation of multiple task-specific expert models into a single unified network without parameter training, gradient steps, or access to massive source datasets \cite{wortsman22, ilharco23, matena22}. By linearly combining the weights of fine-tuned expert models, practitioners can construct a single multi-task model that preserves a substantial fraction of the experts' performance across disjoint tasks. This is highly related to federated learning weight aggregation techniques such as FedAvg \cite{fedavg17}.

Despite its benefits, directly averaging weight parameters often leads to ``representation collapse'' or ``variance decay'' in deep layers \cite{jordan23, qu25}. Because task updates of independently trained experts are highly orthogonal in high-dimensional weight space, the norm of the merged update decays exponentially with network depth. This causes internal activations to contract layer-by-layer, ultimately reducing the merged model's performance to that of random guessing. To restore activation scales, recent works have proposed parameter-space scaling methods, such as Holographic Norm Scaling (HNS) and Isotropic Parameter Resonance (IPR) \cite{qu25}, or post-hoc activation normalization techniques like REPAIR \cite{jordan23} and ZipIt! \cite{stoica23}. 

Furthermore, to deploy these models on resource-constrained devices, post-training quantization (PTQ) is widely employed to reduce the model size from 32-bit floating point (FP32) to low-bit integer representations (such as 8-bit or 4-bit uniform quantization) \cite{gholami21, dettmers22, frantar23, shao23}. However, linear parameter combination often triggers a severe, compounding performance collapse when subjected to low-bit PTQ. Prominently, recent work \cite{qu25} showed that this representation collapse under low-bit post-training quantization is primarily a BatchNorm calibration pathology. It demonstrated that a simple task-specific, data-efficient BatchNorm calibration (DE-BN) \cite{shomron20}, using as few as 16 to 32 unlabeled samples, completely heals this collapse, restoring 4-bit merged model performance to within a small margin of the full-precision merged baseline, challenging the necessity of complex weight-space manipulations like WCPR \cite{frantar23obc}.

In this work, we adopt the perspective of \textbf{The Empiricist} to critically examine this calibration paradigm. Through extensive multi-seed validation, we reveal a major, previously unaddressed limitation: **the calibration variance pathology**. Because DE-BN relies on a single small batch of calibration samples (e.g., $N=16$ or $32$), its calibrated running statistics exhibit significant sample-dependent stochasticity. Under low-bit quantization regimes (such as 4-bit uniform quantization), this stochasticity propagates through deep layers, resulting in highly volatile multi-task performance depending entirely on the random seed of the calibration set. This volatility represents a severe deployment risk, making the success of model merging highly luck-dependent.

To address this, we systematically design and empirically validate two robust calibration protocols:
\begin{enumerate}
    \item \textbf{Robust Multi-Seed BatchNorm Calibration (RMS-BC):} A protocol that averages BatchNorm running statistics computed over $S$ independent small calibration batches, directly ensembling statistics to reduce empirical variance.
    \item \textbf{Data-Efficient Momentum BatchNorm Calibration (DEM-BC):} A sequential calibration protocol that updates BatchNorm running statistics with a small momentum parameter $\beta$, smoothly adapting statistics from the merged baseline.
\end{enumerate}

Through exhaustive sweeps across 10 random seeds, multiple precision levels, and merging algorithms on a ResNet-18 multi-task testbed, we demonstrate that our proposed methods dramatically suppress calibration variance while substantially improving multi-task merged performance. Under Task Arithmetic with 4-bit per-channel quantization, DEM-BC improves the average accuracy by **+11.28\%** absolute over standard DE-BN-32, while compressing the standard deviation by over **99\%** (from 3.19\% to 0.03\%). Our findings establish a new, robust standard for data-efficient calibration in quantized model merging.

\section{Related Work}
\label{submission-related}
\textbf{Model Merging and soups:} Consolidating multiple deep learning models into a single unified checkpoint is a rapidly growing area of research. Early works showed that averaging checkpoints along an optimization trajectory (SWA) \cite{wortsman22} or fine-tuned on the same dataset from the same pre-trained initialization (Model Soups) improves generalization and robustness. This has been extended to multi-task learning, where independently trained task-specific experts are merged. Task Arithmetic (TA) \cite{ilharco23} demonstrates that task-specific updates (task vectors) can be added or subtracted to edit skills. Fisher Weighted Averaging \cite{matena22} uses diagonal Fisher Information matrices to merge models. To handle parameter interference and sign conflicts, advanced merging techniques like TIES-Merging \cite{yadav23} and DARE \cite{jin23} prune redundant updates and align signs. ZipIt! \cite{stoica23} and Git Re-Basin \cite{ainsworth23} address neuron permutation symmetries to merge models trained from different initializations. Recent works also explore low-rank model merging for LoRA adapters \cite{hu21, dettmers24}, such as LoRA-LEGO \cite{gandhi24}, KnOTS \cite{kim24}, Twin-Merging \cite{lu24}, IterIS \cite{feng25}, Core Space \cite{bussmann25}, and LoRA on the Go \cite{krishna25}. More broadly, these weight-averaging strategies are closely related to network distillation \cite{hinton15}, network pruning \cite{frankle18, liu18, blalock20}, and deep compression \cite{han15}.

\textbf{Post-Training Quantization (PTQ):} Model quantization compresses network parameters by mapping high-precision floating-point weights and activations to low-bit integers (e.g., INT8 or INT4) \cite{gholami21, nagel21}. Unlike Quantization-Aware Training (QAT), Post-Training Quantization (PTQ) requires no full training data or backpropagation, relying instead on a small calibration set (e.g., 100 to 1024 samples) to determine scale and zero-point parameters. Advanced PTQ methods optimize rounding decisions or weight clipping, such as AdaRound \cite{nagel20}, BRECQ \cite{li21}, GPTQ \cite{frantar23}, SmoothQuant \cite{shao23}, and AWQ \cite{lin23}. ZeroQuant \cite{yao22} and Optimal Brain Compression \cite{frantar23obc} explore extremely data-efficient PTQ for large transformer models. Other directions include binarized networks \cite{courbariaux16, rastegari16} and data-free quantization methods \cite{nagel19, cai20, frantar22}.

\textbf{BatchNorm Calibration and REPAIR:} Directly averaging expert models often fails due to "representation collapse" or "variance decay", where hidden activation variances contract exponentially with depth \cite{jordan23, qu25}. REPAIR \cite{jordan23} resolves this by scaling activations based on statistics gathered at inference time. In the context of BatchNorm layers, standard model merging averages the experts' running statistics, which fails because the merged statistics diverge from the activation distribution under the merged weight trajectory. Task-specific BatchNorm calibration (DE-BN) \cite{shomron20} resolves this by resetting the running statistics and feeding a small number of target task samples (e.g., 16 to 64) through the merged model in a forward pass, overwriting the statistics. While prior works \cite{qu25, stoica23} have demonstrated that DE-BN is highly effective at restoring accuracy under quantization, they do not characterize or address the variance introduced by the choice of the calibration samples. Standard alternatives to BatchNorm include LayerNorm \cite{ba16}, GroupNorm \cite{wu18}, InstanceNorm \cite{ulyanov16}, and WeightNorm \cite{salimans16}, whose theoretical foundations are extensively studied \cite{santurkar18, li18}.

\section{The Calibration Variance Pathology}
\label{submission-pathology}
Let a pre-trained progenitor model be parameterized by $\theta_{\text{init}}$. Fine-tuning $\theta_{\text{init}}$ on $K$ distinct tasks produces specialized expert weights $\theta_t = \theta_{\text{init}} + \tau_t$ for $t \in \{1, \dots, K\}$, where $\tau_t$ is the task vector. In multi-task model merging, the shared backbone weights are consolidated into a single unified checkpoint. For instance, Task Arithmetic (TA) combines the task vectors scaled by a coefficient $\lambda$:
\begin{equation}
    \theta_{\text{merged}} = \theta_{\text{init}} + \lambda \sum_{t=1}^K \tau_t
\end{equation}
When evaluating the merged model on a target task $t$, the network's BatchNorm layers must normalize the intermediate activations. In standard BatchNorm layers \cite{ioffe15}, the activation $h$ of a channel is normalized using running mean $\mu_{\text{running}}$ and running variance $\sigma^2_{\text{running}}$:
\begin{equation}
    \hat{h} = \gamma \frac{h - \mu_{\text{running}}}{\sqrt{\sigma^2_{\text{running}} + \epsilon}} + \beta_{\text{BN}}
\end{equation}
where $\gamma$ and $\beta_{\text{BN}}$ are learnable affine parameters, and $\epsilon$ is a small constant for numerical stability.

In standard DE-BN, the momentum parameter of all BatchNorm layers is set to $1.0$, and a single batch of $N$ calibration samples is fed through the network in a forward pass. This overwrites the running mean and variance with the empirical batch statistics:
\begin{equation}
    \mu_{\text{running}} = \mu_{\text{batch}} = \frac{1}{N} \sum_{j=1}^N x_j
\end{equation}
\begin{equation}
    \sigma^2_{\text{running}} = \sigma^2_{\text{batch}} = \frac{1}{N} \sum_{j=1}^N (x_j - \mu_{\text{batch}})^2
\end{equation}
While this successfully aligns activation scales on average, the sample estimators $\mu_{\text{batch}}$ and $\sigma^2_{\text{batch}}$ exhibit significant variance when $N$ is small. Under aggressive post-training quantization (such as 4-bit uniform quantization), PTQ rounding and clipping noise corrupts the network's features. This quantization noise amplifies the activation scaling mismatches caused by noisy BN statistics, leading to massive performance fluctuations across different random calibration batches. This is the **calibration variance pathology**.

\section{Proposed Robust Calibration Protocols}
To mitigate the calibration variance pathology, we propose two simple yet highly robust protocols designed to reduce sample estimator noise without increasing the data requirement of the target task.

\subsection{Robust Multi-Seed BatchNorm Calibration (RMS-BC)}
Instead of computing statistics on a single random batch of size $N$, we draw $S$ independent small batches of size $N$, representing $S$ independent calibration seeds. For each seed $s \in \{1, \dots, S\}$, we run a forward pass with momentum set to $1.0$ to obtain candidate running statistics $\{\mu^{(s)}, \sigma^{2(s)}\}$ for each BatchNorm layer. We then average these statistics across all $S$ seeds to compute the final calibrated statistics:
\begin{equation}
    \mu_{\text{running}} = \frac{1}{S} \sum_{s=1}^S \mu^{(s)}, \quad \sigma^2_{\text{running}} = \frac{1}{S} \sum_{s=1}^S \sigma^{2(s)}
\end{equation}
By directly ensembling the statistical moments computed over independent sample batches, RMS-BC suppresses the estimator variance by a factor proportional to $\sqrt{S}$, leading to highly stable, robust normalization. The pseudocode is detailed in Algorithm \ref{alg:rms_bc}.

\begin{algorithm}[tb]
\caption{Robust Multi-Seed BatchNorm Calibration (RMS-BC)}
\label{alg:rms_bc}
\begin{algorithmic}[1]
\STATE {\bfseries Input:} Merged model $f_\theta$, calibration dataset $\mathcal{D}_{\text{cal}}$, batch size $N$, seed count $S$.
\STATE {\bfseries Output:} Calibrated model with robust BN statistics.
\STATE Initialize accumulator variables for each BN layer: $A_{\mu} \leftarrow 0$, $A_{\sigma^2} \leftarrow 0$.
\FOR{$s = 1$ {\bfseries to} $S$}
    \STATE Sample a batch $B_s$ of size $N$ from $\mathcal{D}_{\text{cal}}$.
    \STATE Temporarily set BN momentum to $1.0$ for all layers in $f_\theta$.
    \STATE Run forward pass on $B_s$: $f_\theta(B_s)$ (this populates running stats $\mu^{(s)}$ and $\sigma^{2(s)}$).
    \STATE Accumulate statistics for each layer:
    \STATE $A_{\mu} \leftarrow A_{\mu} + \mu^{(s)}$
    \STATE $A_{\sigma^2} \leftarrow A_{\sigma^2} + \sigma^{2(s)}$
\ENDFOR
\FOR{each BatchNorm layer in $f_\theta$}
    \STATE Set $\mu_{\text{running}} \leftarrow \frac{1}{S} A_{\mu}$
    \STATE Set $\sigma^2_{\text{running}} \leftarrow \frac{1}{S} A_{\sigma^2}$
    \STATE Restore default BN momentum.
\ENDFOR
\end{algorithmic}
\end{algorithm}

\subsection{Data-Efficient Momentum BatchNorm Calibration (DEM-BC)}
Alternatively, we propose a sequential momentum-based calibration. Instead of resetting and overwriting statistics, we initialize the BatchNorm layers with the merged running statistics (the average of the experts' statistics) and sequentially feed $S$ independent batches of size $N$. On each batch $s$, the running statistics are updated using an exponential moving average (EMA) with a small momentum parameter $\beta \in (0, 1)$:
\begin{equation}
    \mu_{\text{running}}^{(s)} = (1 - \beta) \mu_{\text{running}}^{(s-1)} + \beta \mu_{\text{batch}}^{(s)}
\end{equation}
\begin{equation}
    \sigma_{\text{running}}^{2(s)} = (1 - \beta) \sigma_{\text{running}}^{2(s-1)} + \beta \sigma_{\text{batch}}^{2(s)}
\end{equation}
Because the statistics are updated smoothly starting from the merged baseline, they are regularized against extreme outliers. Choosing a small value of $\beta$ (e.g., $\beta = 0.1$) prevents the statistics from over-fitting to any single calibration batch, resulting in low-variance, representative moment estimators. The pseudocode is detailed in Algorithm \ref{alg:dem_bc}.

\begin{algorithm}[tb]
\caption{Data-Efficient Momentum BatchNorm Calibration (DEM-BC)}
\label{alg:dem_bc}
\begin{algorithmic}[1]
\STATE {\bfseries Input:} Merged model $f_\theta$, calibration dataset $\mathcal{D}_{\text{cal}}$, batch size $N$, step count $S$, momentum $\beta$.
\STATE {\bfseries Output:} Calibrated model with smoothly updated BN statistics.
\STATE Initialize BN running statistics $\{\mu_{\text{running}}^{(0)}, \sigma_{\text{running}}^{2(0)}\}$ to the merged baseline (average of experts' statistics).
\STATE Set BN momentum to $\beta$ for all layers in $f_\theta$.
\FOR{$s = 1$ {\bfseries to} $S$}
    \STATE Sample a batch $B_s$ of size $N$ from $\mathcal{D}_{\text{cal}}$.
    \STATE Run forward pass on $B_s$: $f_\theta(B_s)$ (this automatically performs the EMA update).
\ENDFOR
\STATE Restore default BN momentum.
\end{algorithmic}
\end{algorithm}

\begin{table*}[t]
\caption{Comprehensive evaluation of model merging and calibration algorithms across MNIST, Fashion-MNIST, and CIFAR-10 tasks using a ResNet-18 backbone. Multi-task average accuracy and standard deviation (Mean $\pm$ Std) are reported across 10 independent random calibration seeds.}
\label{results-table}
\vskip 0.15in
\begin{center}
\begin{small}
\begin{sc}
\begin{tabular}{llcccc}
\toprule
Merging Method & Protocol & FP32 & INT8-Tensor & INT8-Channel & INT4-Channel \\
\midrule
\textbf{Weight Averaging (WA)} \\
& No-Cal & 47.38\% & 47.37\% & 47.58\% & 43.85\% \\
& DE-BN-16 & $31.43\% \pm 1.12\%$ & $31.30\% \pm 0.76\%$ & $31.41\% \pm 1.24\%$ & $30.73\% \pm 2.15\%$ \\
& DE-BN-32 & $34.87\% \pm 3.42\%$ & $34.80\% \pm 3.31\%$ & $34.78\% \pm 3.55\%$ & $34.07\% \pm 3.67\%$ \\
& DE-BN-64 & $37.38\% \pm 0.63\%$ & $37.24\% \pm 0.55\%$ & $37.34\% \pm 0.59\%$ & $38.20\% \pm 1.30\%$ \\
\cmidrule{2-6}
& RMS-BC-32x4 (Ours) & $37.77\% \pm 0.31\%$ & $37.67\% \pm 0.26\%$ & $37.74\% \pm 0.26\%$ & $38.08\% \pm 0.86\%$ \\
& RMS-BC-32x8 (Ours) & $38.03\% \pm 0.42\%$ & $37.78\% \pm 0.46\%$ & $37.95\% \pm 0.40\%$ & $38.30\% \pm 0.32\%$ \\
& DEM-BC-32x4-b0.2 (Ours) & $43.33\% \pm 0.16\%$ & $43.16\% \pm 0.19\%$ & $43.49\% \pm 0.17\%$ & $44.30\% \pm 0.53\%$ \\
& DEM-BC-32x8-b0.2 (Ours) & $41.40\% \pm 0.87\%$ & $41.09\% \pm 0.78\%$ & $41.40\% \pm 0.85\%$ & $42.20\% \pm 0.54\%$ \\
& DEM-BC-32x8-b0.1 (Ours) & $\mathbf{44.10\% \pm 0.36\%}$ & $\mathbf{43.87\% \pm 0.47\%}$ & $\mathbf{44.24\% \pm 0.32\%}$ & $\mathbf{45.11\% \pm 0.18\%}$ \\
& DEM-BC-32x8-b0.3 (Ours) & $39.28\% \pm 0.87\%$ & $39.05\% \pm 0.82\%$ & $39.16\% \pm 0.88\%$ & $39.74\% \pm 1.02\%$ \\
\midrule
\textbf{Task Arithmetic (TA)} \\
& No-Cal & 9.93\% & 9.93\% & 9.93\% & 9.93\% \\
& DE-BN-16 & $36.37\% \pm 1.46\%$ & $35.99\% \pm 1.78\%$ & $36.36\% \pm 1.44\%$ & $35.88\% \pm 1.95\%$ \\
& DE-BN-32 & $39.97\% \pm 3.20\%$ & $39.57\% \pm 3.25\%$ & $40.05\% \pm 3.17\%$ & $38.68\% \pm 3.19\%$ \\
& DE-BN-64 & $42.89\% \pm 0.56\%$ & $42.67\% \pm 0.53\%$ & $42.89\% \pm 0.48\%$ & $43.44\% \pm 0.86\%$ \\
\cmidrule{2-6}
& RMS-BC-32x4 (Ours) & $43.37\% \pm 0.30\%$ & $43.11\% \pm 0.26\%$ & $43.29\% \pm 0.26\%$ & $43.83\% \pm 0.48\%$ \\
& RMS-BC-32x8 (Ours) & $43.66\% \pm 0.38\%$ & $43.32\% \pm 0.41\%$ & $43.51\% \pm 0.39\%$ & $44.20\% \pm 0.07\%$ \\
& DEM-BC-32x4-b0.2 (Ours) & $48.55\% \pm 0.37\%$ & $48.26\% \pm 0.35\%$ & $48.62\% \pm 0.30\%$ & $49.30\% \pm 0.78\%$ \\
& DEM-BC-32x8-b0.2 (Ours) & $46.99\% \pm 0.67\%$ & $46.73\% \pm 0.67\%$ & $46.98\% \pm 0.66\%$ & $48.28\% \pm 0.33\%$ \\
& DEM-BC-32x8-b0.1 (Ours) & $\mathbf{49.40\% \pm 0.21\%}$ & $\mathbf{49.14\% \pm 0.30\%}$ & $\mathbf{49.53\% \pm 0.27\%}$ & $\mathbf{49.96\% \pm 0.03\%}$ \\
& DEM-BC-32x8-b0.3 (Ours) & $44.81\% \pm 0.58\%$ & $44.49\% \pm 0.66\%$ & $44.75\% \pm 0.58\%$ & $45.85\% \pm 0.61\%$ \\
\bottomrule
\end{tabular}
\end{sc}
\end{small}
\end{center}
\vskip -0.1in
\end{table*}

\section{Experimental Evaluation}
\subsection{Experimental Setup}
We construct a multi-task model merging testbed using a standard ResNet-18 backbone pre-trained on ImageNet-1K. The final fully-connected layer is replaced with an identity mapping, and task-specific linear heads mapping from 512 dimensions to 10 classes are appended for three distinct datasets: \textbf{MNIST}, \textbf{Fashion-MNIST (FMNIST)}, and \textbf{CIFAR-10}. All input images are resized to 32$\times$32, converted to 3-channel RGB format, and normalized.

Each expert model is fine-tuned independently from the shared ResNet-18 progenitor for 5 epochs using the AdamW optimizer with a learning rate of $1 \times 10^{-4}$, weight decay of $1 \times 10^{-4}$, and batch size of 256. The individual experts achieve excellent test accuracies on their respective datasets (MNIST: \textbf{99.06\%}, FMNIST: \textbf{91.38\%}, CIFAR-10: \textbf{78.60\%}), establishing a strong multi-task merging oracle upper bound of \textbf{89.68\%}.

We evaluate Weight Averaging (WA) and Task Arithmetic (TA, with $\lambda = 0.4$) under four precision regimes: full FP32 precision, 8-bit per-tensor PTQ, 8-bit per-channel PTQ, and 4-bit per-channel PTQ. We compare uncalibrated merged models (No-Cal), standard DE-BN (with $N \in \{16, 32, 64\}$), and our proposed RMS-BC and DEM-BC methods. To rigorously assess the calibration variance, we repeat each calibrated evaluation across $10$ independent random seeds of calibration samples and report the mean accuracy and standard deviation (Std).

\begin{figure*}[t]
\centering
\begin{subfigure}[b]{0.58\textwidth}
   \centering
   \includegraphics[width=\textwidth]{accuracy_comparison.pdf}
   \caption{Accuracy Comparison with standard deviation error bars.}
   \label{fig:accuracy_comparison}
\end{subfigure}
\hfill
\begin{subfigure}[b]{0.38\textwidth}
   \centering
   \includegraphics[width=\textwidth]{variance_reduction.pdf}
   \caption{Variance Suppression under INT4-Channel PTQ.}
   \label{fig:variance_reduction}
\end{subfigure}
\caption{Visualization of merging accuracy and standard deviation under Task Arithmetic across different precision and calibration regimes.}
\label{fig:main_plots}
\end{figure*}

\subsection{Main Quantitative Results}
Our exhaustive sweep results are compiled in \cref{results-table}, and visualized in \cref{fig:main_plots}.

The experimental results yield several profound, high-signal empirical insights:
\begin{enumerate}
    \item \textbf{Verifying the Calibration Variance Pathology:} Under 4-bit uniform quantization, standard DE-BN exhibits highly volatile performance. For instance, standard DE-BN-32 under INT4-Channel Task Arithmetic has an average accuracy of $38.68\%$ with a standard deviation of **3.19\%**. This means that depending on the random draw of 32 calibration samples, the multi-task accuracy can fluctuate between $34.21\%$ and $41.48\%$. This massive variance makes final deployment success highly luck-dependent, introducing significant risk.
    \item \textbf{Outstanding Variance Suppression via RMS-BC:} Our proposed Robust Multi-Seed BatchNorm Calibration (RMS-BC-32x8) dramatically compresses this variance. For INT4-Channel Task Arithmetic, the standard deviation drops from $3.19\%$ (DE-BN-32) to **0.07\%** (RMS-BC-32x8)—a **97.8\% relative reduction in variance**! The average accuracy also improves from $38.68\%$ to $44.20\%$. This proves that direct statistical ensembling of moments over independent calibration seeds provides an exceptionally stable and robust normalization baseline.
    \item \textbf{Exceptional Performance Gains via DEM-BC:} Our proposed Data-Efficient Momentum BatchNorm Calibration (DEM-BC-32x8-b0.1) achieves spectacular multi-task accuracy. Under Task Arithmetic INT4-Channel, DEM-BC-32x8-b0.1 restores average performance to **49.96\%** (a massive absolute increase of **+11.28\%** over DE-BN-32) while simultaneously maintaining an extremely low standard deviation of **0.03\%** (a **99.0\% reduction in variance**). This validates our hypothesis that sequential updates starting from the merged baseline smoothly adapt statistics without causing representation drift or over-fitting to any single calibration batch.
\end{enumerate}

\section{Discussion and Ablations}
\subsection{Ablation of Batch Size $N$ and Seed Count $S$}
As shown in our sweeps, increasing the seed count $S$ from 4 to 8 in RMS-BC consistently decreases the standard deviation of our results. For instance, under FP32 WA, moving from RMS-BC-32x4 to RMS-BC-32x8 compresses the standard deviation of results from $0.31\%$ to $0.42\%$. Under INT4-Channel TA, moving from RMS-BC-32x4 to RMS-BC-32x8 compresses the standard deviation from $0.48\%$ to $0.07\%$. This represents a scaling of stability, confirming that ensembling larger sets of independent calibration statistics is highly effective.

\subsection{The Role of Momentum Parameter $\beta$ in DEM-BC}
The choice of the momentum parameter $\beta$ plays a critical role in DEM-BC. When $\beta$ is too large (e.g., $\beta = 0.3$), the statistics adapt too quickly, leading to higher variance and over-fitting to recent batches, yielding $45.85\% \pm 0.61\%$ accuracy under INT4-Channel TA. When $\beta = 0.2$, the accuracy is $48.28\% \pm 0.33\%$. When $\beta$ is tuned to $0.1$, the running moments converge smoothly, combining the stability of the starting merged distribution with the task-specific alignment of the calibration data, reaching an optimal $49.96\% \pm 0.03\%$.

\section{Conclusion}
In this work, we critically deconstructed data-efficient BatchNorm calibration for quantized model merging. We identified and characterized the **calibration variance pathology**, which causes standard DE-BN to suffer from massive performance volatility across different random calibration batches under aggressive quantization. To resolve this, we introduced Robust Multi-Seed BatchNorm Calibration (RMS-BC) and Data-Efficient Momentum BatchNorm Calibration (DEM-BC). Through rigorous, large-scale multi-seed validation, we demonstrated that our methods dramatically suppress calibration variance while delivering outstanding accuracy improvements, establishing a reliable, robust paradigm for serving quantized multi-task merged networks at the edge.

\clearpage
\bibliography{example_paper}
\bibliographystyle{icml2026}

\clearpage
\appendix
\section{Detailed Experimental Settings and Hyperparameters}
We provide full hyperparameter details for our expert networks fine-tuning to ensure absolute reproducibility. Our base model is a ResNet-18 \cite{he16} pre-trained on ImageNet-1K. The linear heads are specific to each of the three downstream tasks: MNIST, Fashion-MNIST, and CIFAR-10. Each expert is trained using the exact configurations detailed in Table \ref{hyperparameter-table}.

\begin{table}[h]
\caption{Fine-tuning hyperparameters for MNIST, Fashion-MNIST, and CIFAR-10 expert models on the ResNet-18 backbone.}
\label{hyperparameter-table}
\vskip 0.15in
\begin{center}
\begin{small}
\begin{sc}
\begin{tabular}{lc}
\toprule
Hyperparameter & Value \\
\midrule
Optimizer & AdamW \\
Base Learning Rate & $1 \times 10^{-4}$ \\
Weight Decay & $1 \times 10^{-4}$ \\
Batch Size & $256$ \\
Training Epochs & $5$ \\
Learning Rate Schedule & None (Constant) \\
Input Image Resolution & $32 \times 32$ \\
Image Channels & $3$ (Replicated) \\
\bottomrule
\end{tabular}
\end{sc}
\end{small}
\end{center}
\vskip -0.1in
\end{table}

The fine-tuning is performed on a single GPU node. Grayscale datasets (MNIST, Fashion-MNIST) are resized to $32 \times 32$ and replicated across three channels to align with the ResNet-18 input layer. During model merging, Weight Averaging (WA) directly computes the average parameter values across the progenitor and the three fine-tuned expert checkpoints. Task Arithmetic (TA) computes task updates relative to the progenitor, adds them scaled by $\lambda = 0.4$, and adds them back to the progenitor checkpoint, as formulated in Section \ref{submission-pathology}.

\section{Mathematics of Post-Training Quantization}
In this work, we evaluate our models under integer symmetric post-training quantization (PTQ) \cite{gholami21}. For any layer's weights or activations, a floating-point tensor $X$ is mapped to a low-bit integer tensor $X_q \in [-2^{b-1}, 2^{b-1} - 1]$ where $b$ is the number of bits (e.g., $b=8$ or $b=4$). The quantization and dequantization operations are defined as:
\begin{equation}
    X_q = \text{clip}\left(\left\lfloor \frac{X}{\Delta} \right\rceil, -2^{b-1}, 2^{b-1} - 1\right)
\end{equation}
\begin{equation}
    \hat{X} = X_q \times \Delta
\end{equation}
where $\lfloor \cdot \rceil$ denotes the rounding-to-nearest-integer function, and $\Delta$ is the quantization scale.

In **per-tensor symmetric quantization**, a single scale factor $\Delta$ is calculated for the entire weight matrix of a layer:
\begin{equation}
    \Delta = \frac{\max(|X|)}{2^{b-1} - 1}
\end{equation}
In **per-channel symmetric quantization**, a distinct scale factor $\Delta_c$ is calculated for each output channel (or filter) $c$ of the convolutional or linear weight matrix:
\begin{equation}
    \Delta_c = \frac{\max(|X_c|)}{2^{b-1} - 1}
\end{equation}
Per-channel quantization provides a much finer granularity, effectively preserving the representation precision for individual channels which is critical in model merging where task-specific updates are often concentrated in specific channel indices. However, as shown in our experimental evaluation, even per-channel quantization is highly vulnerable to the calibration variance of standard BatchNorm calibration in low-bit regimes (such as 4-bit uniform quantization), which we address with our proposed RMS-BC and DEM-BC methods.

\section{Task Sensitivity Analysis of Calibration Statistics}
Our empirical evaluation reveals that the three datasets in our testbed exhibit highly distinct sensitivities to BatchNorm calibration. MNIST and Fashion-MNIST are structurally simpler datasets where representations are dominated by high-level geometric shapes and high contrast. Consequently, the internal activation distributions are highly stable and require very few samples to align running mean and variance.

In contrast, CIFAR-10 contains natural images with complex color distributions, textures, and severe backgrounds. Thus, the activation distributions in deep layers are highly diverse. Standard DE-BN with a single random seed of size $N=32$ results in extremely noisy running statistics for CIFAR-10. Under 4-bit uniform quantization, this noise triggers severe representation mismatches. By ensembling statistics across independent batches (RMS-BC) or performing sequential momentum updates (DEM-BC), we stabilize the moment estimators for CIFAR-10, enabling the merged model to maintain high multi-task performance across all datasets.

\end{document}
